\testcase{T01 --- Wyszukanie pasującego auta}
{Dostępne jest auto zgodne z kryteriami.}
{BuyerAgent, SellerAgent.}
{Kupujący wysyła \code{car-search}.}
{Sprzedający zwraca PROPOSE z najtańszą zgodną ofertą.}
{Oferta została zwrócona i użyta w rankingu.}{PASS}

\testcase{T02 --- Brak pasującego auta}
{Żadne dostępne auto nie spełnia filtrów.}
{BuyerAgent, SellerAgent.}
{Kupujący wysyła zapytanie z niemożliwymi kryteriami.}
{Odpowiedź REFUSE / NO\_OFFER.}
{Brak oferty obsłużono bez przerwania symulacji.}{PASS}

\testcase{T03 --- Zakup w ramach budżetu}
{Saldo przekracza uzgodniony koszt całkowity.}
{BuyerAgent, SellerAgent.}
{Rezerwacja, zakup i płatność.}
{Auto otrzymuje SOLD, a saldo maleje o koszt transakcji.}
{Transakcja została utworzona; saldo jest poprawne.}{PASS}

\testcase{T04 --- Zakup przekraczający budżet}
{Koszt całkowity przekracza saldo.}
{\code{BuyerAccount}, komponent budżetu kupującego.}
{Próba płatności ponad dostępne środki.}
{Odrzucenie operacji bez zmiany salda.}
{\code{IllegalStateException}; saldo pozostało bez zmian.}{PASS}

\testcase{T05 --- Konkurencja o jeden samochód}
{Kilku kupujących wybiera tę samą ofertę.}
{Kilka BuyerAgent i jeden SellerAgent.}
{Równoległe żądania rezerwacji.}
{Tylko jedno AGREE; pozostałe żądania otrzymują REFUSE.}
{W logu odnotowano odrzucone konkurencyjne rezerwacje.}{PASS}

\testcase{T06 --- Rezerwacja samochodu}
{Auto AVAILABLE, poprawna uzgodniona cena.}
{BuyerAgent, SellerAgent.}
{Wysłanie \code{car-reservation}.}
{Status RESERVED i przypisanie właściciela rezerwacji.}
{Rezerwacja została zaakceptowana.}{PASS}

\testcase{T07 --- Wygaśnięcie rezerwacji}
{Auto zarezerwowane bez późniejszego zakupu.}
{\code{SellerInventory}, komponent sprzedającego.}
{Przesunięcie zegara poza pięciosekundowy okres rezerwacji.}
{Automatyczny powrót do AVAILABLE.}
{Test sterowanego zegara potwierdził zwolnienie pojazdu.}{PASS}

\testcase{T08 --- Udana negocjacja}
{Możliwe porozumienie w granicach maksymalnego rabatu 15\%.}
{BuyerAgent, NegotiatorAgent; test reguł \code{NegotiationService}.}
{Oferta kupującego i kontrpropozycja.}
{Akceptacja w najwyżej trzech rundach.}
{Cena została zaakceptowana i użyta w transakcji. Test usługi
potwierdza kontrpropozycję 133\,500 PLN przy pierwszej ofercie
119\,000 PLN i limicie 148\,000 PLN.}{PASS}

\testcase{T09 --- Odrzucona negocjacja}
{Limit kupującego jest niższy niż minimalna dopuszczalna cena.}
{BuyerAgent, NegotiatorAgent; \code{NegotiationService}.}
{Negocjacja z pierwszą ofertą 110\,000 PLN i limitem 120\,000 PLN.}
{Odrzucenie najpóźniej po trzech rundach.}
{\code{accepted=false}; trzy rundy i brak uzgodnionej ceny.}{PASS}

\testcase{T10 --- Rekomendacja TOP 3}
{Co najmniej cztery oferty SUV spełniają kryteria.}
{BuyerAgent, AdvisorAgent; \code{RecommendationService}.}
{Przekazanie ofert i preferencji doradcy.}
{Trzy oferty o najwyższej ocenie MCDA w poprawnej kolejności.}
{Zwrócono CAR-001, CAR-002 i CAR-003; niezgodny SEDAN odrzucono.}{PASS}

\testcase{T11 --- Ocena sprzedawcy}
{Potwierdzone transakcje umożliwiają wystawienie ocen.}
{BuyerAgent, ReputationAgent; \code{ReputationService}.}
{Oceny 5 i 4 dla sprzedawcy oraz ponowna ocena pierwszej transakcji.}
{Średnia 4,5, dwie oceny i odrzucenie duplikatu.}
{Średnia oraz licznik są poprawne; duplikat odrzucono.}{PASS}

\testcase{T12 --- Analiza trendu}
{Cztery sprzedaże o rosnących cenach.}
{BuyerAgent, TrendAgent; \code{TrendService}.}
{Kolejne zdarzenia sprzedaży: 100\,000, 105\,000, 130\,000 i 140\,000 PLN.}
{Poprawne popularne kategorie, średnia i trend RISING.}
{Średnia wyniosła 118\,750 PLN; najpopularniejsze są SUV,
Toyota RAV4 i PETROL; trend RISING.}{PASS}

\testcase{T13 --- Zwrot w terminie}
{Od zakupu nie upłynęło pięć sekund.}
{BuyerAgent, SellerAgent.}
{Właściciel wysyła \code{car-return}.}
{Zwrot środków, AVAILABLE i historia RETURNED.}
{Wykonano pełny zwrot; kupujący następnie kupił zamiennik.}{PASS}

\testcase{T14 --- Zwrot po terminie}
{Od zakupu minęło ponad pięć sekund.}
{BuyerAgent, SellerAgent.}
{Żądanie zwrotu po upływie okna czasowego.}
{Odmowa i pozostawienie auta w stanie SOLD.}
{Otrzymano \code{RETURN_WINDOW_EXPIRED}.}{PASS}

\testcase{T15 --- Określenie wartości rynkowej}
{Oferta oraz dwa porównywalne pojazdy.}
{AdvisorAgent; \code{ValuationService}.}
{Obliczenie średniej i procentowego odchylenia ceny.}
{Poprawna kwalifikacja do przedziałów 0--5, 5--15 i powyżej 15\%.}
{Zweryfikowano granice przedziałów klasyfikacji.}{PASS}

\testcase{T16 --- Wybór metodą MCDA}
{Oferty różnią się ceną, rokiem, paliwem, pojemnością i spalaniem.}
{BuyerAgent, AdvisorAgent; \code{RecommendationService}.}
{Normalizacja cech i obliczenie ważonej sumy.}
{Ranking zgodny z wagami 30/20/15/10/10/15.}
{Kolejność rekomendacji i zakres ocen 0--1 są poprawne.}{PASS}

\testcase{T17 --- Jednoczesna komunikacja wielu agentów}
{10 sprzedających, 5 kupujących i 80 samochodów.}
{Wszystkie agenty aplikacyjne.}
{Uruchomienie pełnej symulacji.}
{15 zachowanych zakupów, bez podwójnej sprzedaży i pozostałych rezerwacji.}
{15 SOLD, 0 RESERVED; każdy kupujący ma trzy auta.}{PASS}

\subsection{Podsumowanie testów}
Dokumentacja źródłowa podaje 17 scenariuszy zaliczonych i zero niezaliczonych.
Scenariusze obejmują zarówno ścieżki poprawne, jak i odmowę z powodu
braku ofert, niewystarczającego budżetu, konkurencji, niedopuszczalnej
ceny lub przekroczenia czasu zwrotu.
